% BHU PYQ -- Economics: Indian Economy (Major-2) (2025-26 NEP)
\documentclass[11pt,a4paper]{article}
\usepackage[a4paper,top=2.5cm,bottom=2.5cm,left=2.8cm,right=2.8cm,headheight=14pt]{geometry}
\usepackage{amsmath,amssymb}
\usepackage{fontspec}
\setmainfont{Kohinoor Devanagari}
\usepackage{microtype,fancyhdr,enumitem,hyperref,lastpage,parskip}

\hypersetup{colorlinks=true,urlcolor=black,linkcolor=black,pdftitle={ECOMJ32: Indian Economy (Major-2) -- BHU NEP 2025-26}}

\pagestyle{fancy}\fancyhf{}
\renewcommand{\headrulewidth}{0.4pt}\renewcommand{\footrulewidth}{0.4pt}
\fancyhead[L]{\small \textit{Banaras Hindu University}}
\fancyhead[R]{\small \textit{ECOMJ32: Indian Economy (2025-26)}}
\fancyfoot[C]{\small Page \thepage\ of \pageref{LastPage}}

\newlist{parts}{enumerate}{2}
\setlist[parts,1]{label=(\alph*),leftmargin=*,itemsep=4pt,topsep=3pt}
\setlist[parts,2]{label=(\roman*),leftmargin=*,itemsep=2pt,topsep=2pt}

\newcommand{\pts}[1]{\hfill \textit{[#1]}}

\begin{document}

\begin{center}
  {\large\bfseries BANARAS HINDU UNIVERSITY}\\[2pt]
  {\normalsize B. A. (Hons.) Semester III Examination, 2025-26 (NEP)}\\[6pt]
  \rule{0.8\linewidth}{0.5pt}\\[6pt]
  {\large\bfseries ECONOMICS}\\[4pt]
  {\large\bfseries Paper Code: ECOMJ32 : Indian Economy (Major-2)}\\[4pt]
  {\normalsize Time: 2:30 Hours \hfill Full Marks: 70}\\[2pt]
  {\small REF NO. 2025-26/D-III/032}
\end{center}
\rule{\linewidth}{0.4pt}
\smallskip

\noindent \textit{Instructions: This question paper comprises three Sections A, B and C. All questions are compulsory. Write your Roll No.\ at the top immediately on receipt of this question paper.}

\smallskip
\noindent \textit{इस प्रश्न-पत्र में तीन खण्ड अ, ब तथा स हैं। सभी प्रश्न अनिवार्य हैं।}

\bigskip

\section*{SECTION -- A / खण्ड -- अ}
\noindent \textit{Note: Answer each question within 50 words. Each question carries 2 marks.}\pts{10 $\times$ 2 = 20}

\noindent \textit{प्रत्येक प्रश्न का उत्तर 50 शब्दों के भीतर दीजिए। प्रत्येक प्रश्न 2 अंकों का है।}

\begin{enumerate}[label=\textbf{\arabic*.},leftmargin=*,itemsep=8pt]
  \item Explain the following:\pts{10 $\times$ 2 = 20}
  \begin{parts}
    \item White Revolution / श्वेत क्रान्ति
    \item 7 Pillars of NITI Aayog / नीति आयोग के 7 स्तम्भ
    \item Disinvestment / विनिवेश
    \item Blue Revolution / नीली क्रान्ति
    \item Foreign Portfolio Investment (FPI) / विदेशी पोर्टफोलियो निवेश
    \item Balance of Payments / भुगतान संतुलन
    \item Agricultural Crisis / कृषि संकट
    \item Demonetization / विमुद्रीकरण
    \item Merchandise Trade / व्यापारिक माल का व्यापार
    \item Foreign Direct Investment (FDI) / प्रत्यक्ष विदेशी निवेश
  \end{parts}
\end{enumerate}

\bigskip

\section*{SECTION -- B / खण्ड -- ब}
\noindent \textit{Note: Write your answer in about 125 words. Each question carries 5 marks.}\pts{4 $\times$ 5 = 20}

\noindent \textit{अपना उत्तर लगभग 125 शब्दों में लिखिए। प्रत्येक प्रश्न 5 अंकों का है।}

\begin{enumerate}[label=\textbf{\arabic*.},leftmargin=*,start=2,itemsep=12pt]

  \item Discuss five major negative impacts of MNCs on the Indian Economy.\pts{5}
  
  भारतीय अर्थव्यवस्था पर बहुराष्ट्रीय कंपनियों (MNCs) के पाँच प्रमुख नकारात्मक प्रभावों की चर्चा कीजिए।

  \begin{center}\textbf{OR / अथवा}\end{center}

  Discuss the features of the Indian Economy on the eve of independence.\pts{5}
  
  स्वतंत्रता की पूर्व संध्या पर भारतीय अर्थव्यवस्था की विशेषताओं की चर्चा कीजिए।

  \item ``Indian economy has leapfrogged from the agricultural sector directly to the tertiary sector.'' Discuss.\pts{5}
  
  ``भारतीय अर्थव्यवस्था कृषि क्षेत्र से सीधे तृतीयक (सेवा) क्षेत्र में छलांग लगा चुकी है।'' विवेचना कीजिए।

  \begin{center}\textbf{OR / अथवा}\end{center}

  Discuss the major causes of low productivity in Indian agriculture.\pts{5}
  
  भारतीय कृषि में निम्न उत्पादकता के प्रमुख कारणों की चर्चा कीजिए।

  \item What do you mean by Black Economy? Discuss the major causes of Black Economy in India.\pts{5}
  
  काली अर्थव्यवस्था (काला धन) से आप क्या समझते हैं? भारत में काली अर्थव्यवस्था के प्रमुख कारणों की चर्चा कीजिए।

  \begin{center}\textbf{OR / अथवा}\end{center}

  Discuss the first three Five-Year Plan models of India.\pts{5}
  
  भारत की प्रथम तीन पंचवर्षीय योजनाओं के मॉडलों की चर्चा कीजिए।

  \item Define Public Sector. Discuss the main problems of Public Sector Enterprises in India.\pts{5}
  
  सार्वजनिक क्षेत्र को परिभाषित कीजिए। भारत में सार्वजनिक क्षेत्र के उपक्रमों की मुख्य समस्याओं की चर्चा कीजिए।

\end{enumerate}

\bigskip

\section*{SECTION -- C / खण्ड -- स}
\noindent \textit{Note: Answer the following questions in about 500 words each. Each question carries 10 marks.}\pts{3 $\times$ 10 = 30}

\noindent \textit{निम्नलिखित प्रश्नों के उत्तर लगभग 500 शब्दों में दीजिए। प्रत्येक प्रश्न 10 अंकों का है।}

\begin{enumerate}[label=\textbf{\arabic*.},leftmargin=*,start=6,itemsep=14pt]

  \item Critically evaluate the measures taken to increase crop production and productivity under the Green Revolution in India.\pts{10}
  
  भारत में हरित क्रान्ति के अंतर्गत फसलों के उत्पादन एवं उत्पादकता को बढ़ाने के लिए उठाए गए उपायों का समालोचनात्मक मूल्यांकन कीजिए।

  \begin{center}\textbf{OR / अथवा}\end{center}

  Discuss the major economic reforms of 1991 (LPG: Liberalization, Privatization, and Globalization) and their impact on India's growth rate.\pts{10}
  
  1991 के प्रमुख आर्थिक सुधारों (उदारीकरण, निजीकरण एवं वैश्वीकरण) तथा भारत की संवृद्धि दर पर उनके प्रभाव की विवेचना कीजिए।

  \item Discuss the trends and sectoral composition of National Income in India since independence.\pts{10}
  
  स्वतंत्रता के बाद से भारत में राष्ट्रीय आय की प्रवृत्तियों तथा क्षेत्रीय संरचना (Sectoral Composition) की विवेचना कीजिए।

  \begin{center}\textbf{OR / अथवा}\end{center}

  Examine the problems of Poverty and Unemployment in India. Discuss the major government welfare schemes aimed at poverty alleviation.\pts{10}
  
  भारत में गरीबी एवं बेरोजगारी की समस्याओं का परीक्षण कीजिए। गरीबी उन्मूलन के उद्देश्य से संचालित प्रमुख सरकारी कल्याणकारी योजनाओं की चर्चा कीजिए।

\end{enumerate}

\end{document}
